\section{Related Work}\label{sec:related-work}

The study of asynchronous programming sits at the intersection of several bodies of literature. We organize related work into three groups: foundational models of concurrency from which asynchrony draws its semantics, the development of deferred-value abstractions that gave asynchrony its first concrete programming primitives, and the functional and monadic tradition that supplied the compositional framework underlying async/await.

The earliest formal treatments of concurrency were not primarily concerned with asynchrony as we have defined it, but they established the semantic vocabulary on which later work depends. The first formalized models were the actor model~\cite{baker_incremental_1977,agha_actors_1986} and communicating sequential processes~\cite{hoare_communicating_1978}. The $\pi$-calculus extended this work by allowing first-class channels that can also be sent between processes~\cite{milner_calculus_1992}.

Deferred value abstractions have been introduced in many forms and under many aliases. \citet{baker_incremental_1977} introduced the term \textit{future,} which was also used by Mulitlisp, that extended Scheme with several forms for expressing concurrency~\cite{halstead_multilisp_1985}. An independent exploration of demand-driven evaluation by \missingcite brought about the lazy \textit{cons} operator. \citet{hutchison_concurrency_2005} defined \textit{promise pipelining,} rather than blocking to receive a value a program could register a continuation that would execute upon resolution.

In the functional world rather than asking how to represent a deferred value, the exploration asked how to compose effectful computations while maintaining referential transparency. \citet{jones_concurrent_1996} built concurrent Haskell, which was later extended with shared-state concurrency~\cite{harris_composable_nodate}. The monadic model influenced \FSharp, the first major async/await design, where the compiler automatically rewrote a block of code written in direct-style into its continuation-passing equivalent~\cite{syme_f_2011}. \citet{bierman_pause_2012} published a formal model for featherweight \CSharp$_{5}$ from which we drew much inspiration.

